% !TeX program = xelatex
% 编译两次以更新目录及总页数。UTF-8，推荐 TeX Live。
\documentclass[aspectratio=169,11pt,notheorems]{beamer}
\usepackage[UTF8,fontset=fandol]{ctex}
\usepackage{booktabs,listings,tikz,algpseudocode}
\usetikzlibrary{arrows.meta,positioning}
\usefonttheme{professionalfonts}

% ── 配色与画布 ─────────────────────────────
\definecolor{Paper}{HTML}{F7F7F7}
\definecolor{Ink}{HTML}{23373D}
\definecolor{Navy}{HTML}{1D486F}
\definecolor{Accent}{HTML}{CB4542}
\definecolor{Forest}{HTML}{34703B}
\definecolor{Track}{HTML}{CED8DE}
\setbeamercolor{normal text}{fg=Ink,bg=Paper}
\setbeamercolor{structure}{fg=Navy}
\setbeamercolor{frametitle}{fg=white,bg=Navy}
\setbeamercolor{progress}{fg=Accent,bg=Track}
\setbeamercolor{footline}{fg=Ink,bg=Paper}
\setbeamerfont{frametitle}{size=\large,series=\bfseries}
\setbeamersize{text margin left=10mm,text margin right=10mm}
\setbeamertemplate{navigation symbols}{}
\setbeamertemplate{headline}{}
\setbeamertemplate{section in toc}{%
  \leavevmode\makebox[2.3em][l]{\color{Navy}\bfseries\inserttocsectionnumber}%
  \inserttocsection\par\vspace{.5em}}
\setbeamertemplate{itemize item}{\raise.2ex\hbox{\color{Ink}\rule{.7mm}{.7mm}}}
\setbeamertemplate{itemize subitem}{\color{Navy}\textendash}
\setbeamertemplate{enumerate item}{\color{Navy}\bfseries\insertenumlabel.}
\setbeamertemplate{enumerate subitem}{\color{Navy}\insertsubenumlabel.}

% ── 页眉、进度线、页码 ─────────────────────
\newcommand{\SlideProgress}{%
  \pgfmathsetlengthmacro{\DoneWidth}{%
    \paperwidth*min(1,\insertframenumber/max(1,\inserttotalframenumber))}%
  \begin{beamercolorbox}[wd=\paperwidth,ht=.35pt,dp=0pt]{progress}%
    \color{Accent}\rule{\DoneWidth}{.35pt}%
  \end{beamercolorbox}}
\setbeamertemplate{frametitle}{%
  \nointerlineskip
  \begin{beamercolorbox}[wd=\paperwidth,sep=2.3mm,leftskip=1.1mm]{frametitle}%
    \usebeamerfont{frametitle}\insertframetitle\par
    \ifx\insertframesubtitle\empty\else
      {\footnotesize\mdseries\insertframesubtitle\par}%
    \fi
  \end{beamercolorbox}%
  \nointerlineskip\SlideProgress}
\setbeamertemplate{footline}{%
  \begin{beamercolorbox}[wd=\paperwidth,ht=3mm,dp=3mm,rightskip=2.5mm plus 1fil]{footline}%
    \hfill\scriptsize\insertframenumber/\inserttotalframenumber
  \end{beamercolorbox}}

% ── 封面与分节页 ───────────────────────────
\setbeamertemplate{title page}{%
  {\Large\bfseries\inserttitle\par}
  \vspace{.55em}{\large\insertsubtitle\par}
  \vspace{1em}{\color{Accent}\rule{\linewidth}{.45pt}\par}
  \vspace{1.2em}{\small\insertauthor\par}
  \vspace{.65em}{\footnotesize\insertinstitute\par\insertdate\par}}
\newcommand{\ChapterPage}[1]{%
  \begin{frame}[plain,noframenumbering]
    \centering
    \begin{minipage}{.62\textwidth}
      {\large\bfseries #1\par}\vspace{.5em}
      {\color{Accent}\rule{2.5mm}{.45pt}}%
      {\color{Track}\rule{\dimexpr\linewidth-2.5mm\relax}{.45pt}}
    \end{minipage}
  \end{frame}}
\AtBeginSection[]{\ChapterPage{第 \insertsectionnumber\ 部分：\insertsectionhead}}

% ── 块、代码与图形 ─────────────────────────
\setbeamercolor{block title}{fg=Navy,bg=Navy!12!Paper}
\setbeamercolor{block body}{fg=Ink,bg=Navy!5!Paper}
\setbeamercolor{block title example}{fg=Forest,bg=Forest!13!Paper}
\setbeamercolor{block body example}{fg=Ink,bg=Forest!6!Paper}
\setbeamercolor{block title alerted}{fg=Accent,bg=Accent!13!Paper}
\setbeamercolor{block body alerted}{fg=Ink,bg=Accent!6!Paper}
\newcommand{\PanelStart}[1]{%
  \par\medskip
  \begin{beamercolorbox}[sep=1.2mm,leftskip=1mm]{block title#1}%
    \bfseries\insertblocktitle
  \end{beamercolorbox}%
  \nointerlineskip
  \begin{beamercolorbox}[sep=1.4mm,leftskip=.8mm,rightskip=.8mm]{block body#1}}
\newcommand{\PanelEnd}{\end{beamercolorbox}\par\smallskip}
\setbeamertemplate{block begin}{\PanelStart{}}
\setbeamertemplate{block end}{\PanelEnd}
\setbeamertemplate{block example begin}{\PanelStart{ example}}
\setbeamertemplate{block example end}{\PanelEnd}
\setbeamertemplate{block alerted begin}{\PanelStart{ alerted}}
\setbeamertemplate{block alerted end}{\PanelEnd}
\newtheorem{theorem}{定理}[section]
\newtheorem{definition}[theorem]{定义}
\renewcommand{\proofname}{证明}
\setbeamercolor{alerted text}{fg=Accent}
\lstdefinestyle{lecture}{
  basicstyle=\ttfamily\small, backgroundcolor=\color{Navy!5!Paper},
  keywordstyle=\color{Navy}\bfseries, commentstyle=\color{Forest},
  stringstyle=\color{Accent}, numbers=left, numberstyle=\tiny\color{Navy},
  xleftmargin=6mm, framexleftmargin=5mm, frame=single,
  rulecolor=\color{Track}, showstringspaces=false, breaklines=true, columns=fullflexible
}
\tikzset{
  lecture step/.style={draw=Navy,fill=Navy!6!Paper,rounded corners=1.3mm,
    align=center,minimum height=10mm,text width=22mm,font=\small},
  lecture focus/.style={lecture step,draw=Accent,fill=Accent!7!Paper},
  lecture flow/.style={-{Stealth[length=2mm]},draw=Navy,semithick}
}

% ── 替换以下信息与示例内容 ─────────────────
\title{讲义标题：从问题到实现}
\subtitle{副标题或一句话摘要}
\author{主讲人：姓名}
\institute{单位／课程名称}
\date{\today}

\begin{document}
\begin{frame}[plain,noframenumbering]
  \titlepage
\end{frame}

\begin{frame}{目录}
  \tableofcontents
\end{frame}

\section{问题与记号}

\begin{frame}{背景与目标}{可选的小标题}
  \begin{itemize}
    \item 输入：一组观测及其特征表示。
    \item 输出：每个样本所属的类别。
      \begin{itemize}
        \item 根据距离度量比较样本。
        \item 用统一指标评估结果。
      \end{itemize}
    \item \alert{注意：}正式比较之前，统一预处理和评价条件。
  \end{itemize}
\end{frame}

\begin{frame}{定义、定理与证明}
  \begin{definition}
    若 $\sum_{i=1}^{n}x_i^2=1$，则称 $x$ 为单位向量。
  \end{definition}
  \begin{theorem}[单位向量的内积]
    对任意单位向量 $x,y$，有 $|x^\top y|\leq 1$。
  \end{theorem}
  \begin{proof}
    由柯西–施瓦茨不等式，$|x^\top y|\leq\|x\|_2\|y\|_2=1$。
  \end{proof}
\end{frame}

\begin{frame}{说明块与枚举}
  \begin{exampleblock}{示例}
    向量 $(3,4)$ 归一化后得到 $(0.6,0.8)$。
  \end{exampleblock}
  \begin{alertblock}{边界条件}
    零向量应单独处理，以避免除零。
  \end{alertblock}
  \begin{enumerate}
    \item 检查输入数据。
    \item 计算范数并归一化。
    \item 验证输出并记录异常。
  \end{enumerate}
\end{frame}

\section{流程与实现}

\begin{frame}{处理流程}
  \centering
  \begin{tikzpicture}[node distance=5mm]
    \node[lecture step] (input) {输入\\原始数据};
    \node[lecture step,right=of input] (features) {预处理\\特征表示};
    \node[lecture focus,right=of features] (method) {核心算法\\计算结果};
    \node[lecture step,right=of method] (output) {输出\\结果与指标};
    \draw[lecture flow] (input) -- (features);
    \draw[lecture flow] (features) -- (method);
    \draw[lecture flow] (method) -- (output);
  \end{tikzpicture}
  \par\bigskip
  {\small\color{Navy}蓝色：通用步骤\qquad\textcolor{Accent}{红色：核心步骤}}
\end{frame}

% 含 lstlisting 的 frame 要保留 [fragile]。
\begin{frame}[fragile]{代码示例：向量归一化}
\begin{lstlisting}[style=lecture,language=Python]
from math import sqrt

def normalize(values):
    norm = sqrt(sum(x * x for x in values))
    if norm == 0:
        raise ValueError("zero vector")
    return [x / norm for x in values]
\end{lstlisting}
\end{frame}

\begin{frame}{算法：向量归一化}
  \begin{block}{算法 1\quad Normalize}
    \begin{algorithmic}[1]
      \Require 非零向量 $x=(x_1,\ldots,x_n)$
      \Ensure 单位向量 $y$
      \State $r\gets\sqrt{\sum_{i=1}^{n} x_i^2}$
      \For{$i=1,\ldots,n$}
        \State $y_i\gets x_i/r$
      \EndFor
      \State \Return $y$
    \end{algorithmic}
  \end{block}
\end{frame}

\section{结果与总结}

\begin{frame}{实验结果}
  \centering
  \renewcommand{\arraystretch}{1.32}
  \begin{tabular}{@{}lrrr@{}}
    \toprule
    \textbf{方法} & \textbf{准确率} & \textbf{召回率} & \textbf{用时（秒）} \\
    \midrule
    基线方法 & 81.2\% & 79.4\% & 1.24 \\
    改进方法 & \textbf{86.7\%} & \textbf{85.1\%} & 1.38 \\
    \bottomrule
  \end{tabular}
  \par\medskip{\footnotesize 数值仅用于演示表格样式。}
\end{frame}

\begin{frame}{总结}
  \begin{itemize}
    \item 明确输入、输出和边界条件。
    \item 用流程图解释方法，用伪代码说明实现。
    \item 在统一条件下比较实验结果。
  \end{itemize}
\end{frame}

\ChapterPage{谢谢}
\end{document}
